% MATH Companion Part VI / VI-32 / explanatory bridge only.
\paragraph{Why differentials belong in a chapter about tangent spaces.}
A tangent vector describes an infinitesimal first-order change of
a function.  Algebraically such changes satisfy the Leibniz rule.
For a field \(k\), a \(k\)-derivation
\(D:B\to M\), with \(M\) a \(B\)-module, is \(k\)-linear
and satisfies
\[
 D(fg)=fD(g)+gD(f).
\]
The module of K\"ahler differentials \(\Omega_{B/k}\) is
characterized by its universal property:
\[
 \operatorname{Der}_k(B,M)
 \cong\operatorname{Hom}_B(\Omega_{B/k},M).
\]
Thus differential forms and derivations are dual ways of encoding
first-order information.  This is the reason the construction
\(I/I^2\), where \(I=\ker(B\otimes_k B\to B)\), appears naturally
in the next proof.

\paragraph{The cotangent space at a rational point.}
Let \(x\) be a \(k\)-rational point of an affine
\(k\)-scheme \(X=\operatorname{Spec}B\), with maximal ideal
\(\mathfrak m\) and residue field \(B/\mathfrak m=k\).
The cotangent space is \(\mathfrak m/\mathfrak m^2\).
Its dual vector space
\[
 T_xX=\operatorname{Hom}_k
 (\mathfrak m/\mathfrak m^2,k)
\]
is the \emph{Zariski tangent space}.  Modding out by
\(\mathfrak m^2\) discards second-order and higher terms:
only the linear part of a function near \(x\) remains.
At a rational point there is a canonical identification
\[
 \Omega_{B/k}\otimes_B k\cong\mathfrak m/\mathfrak m^2.
\]
For non-rational points one must keep track of the residue field
and the exact sequence of cotangent spaces; the displayed
simplification is not asserted without its rational-point hypothesis.

\paragraph{Deriving the Jacobian rule.}
For a quotient \(B=k[x_1,\ldots,x_n]/(f_1,\ldots,f_r)\),
the relations among the free differentials on affine \(n\)-space
are obtained by differentiating the defining equations:
\[
 \Omega_{B/k}\cong
 \frac{\bigoplus_{i=1}^n B\,dx_i}
 {\langle df_1,\ldots,df_r\rangle},
 \qquad
 df_j=\sum_{i=1}^{n}
 \frac{\partial f_j}{\partial x_i}\,dx_i.
\]
Tensoring with the residue field at a rational point \(a\)
replaces each coefficient by its value at \(a\).  The linear
relations are the rows of the Jacobian matrix
\(J(a)=(\partial f_j/\partial x_i)(a)\), so
\[
 \boxed{\dim_k T_aX=n-\operatorname{rank}J(a).}
\]
This formula describes the tangent space of a specific affine
presentation.  For a finite-type scheme over a characteristic-zero
field, tangent-space dimension is at least the local dimension at a
closed point; equality is the expected regularity test.

\paragraph{Worked comparison: a smooth point and a cusp.}
Assume \(\operatorname{char}k=0\), and let
\(C=V(y^2-x^3)\subseteq\mathbb A^2_k\).  Writing
\(f=y^2-x^3\) gives
\[
 J(x,y)=\begin{pmatrix}-3x^2&2y\end{pmatrix}.
\]
At \((1,1)\), the Jacobian has rank one.  The tangent
space is therefore one-dimensional, in agreement with the local
dimension of the curve.  At the singular point \((0,0)\), both
partial derivatives vanish, so the Jacobian has rank zero and
\[
 \dim_k T_{(0,0)}C=2,
 \qquad\dim\mathcal O_{C,(0,0)}=1.
\]
The curve has not become two-dimensional.  Rather, the first-order
equation vanishes at the cusp, so the tangent-space approximation
is larger than the actual local geometry.  This is a concrete
reason to keep Krull dimension and embedding dimension separate.

\paragraph{Seeing the same phenomenon through differentials.}
The relation \(y^2-x^3=0\) yields
\[
 \Omega_{B/k}\cong
 \frac{B\,dx\oplus B\,dy}
 {\langle 2y\,dy-3x^2\,dx\rangle}.
\]
At the cusp, tensoring with \(B/(x,y)=k\) makes the coefficient
of this relation zero.  The classes of \(dx\) and \(dy\) then
remain linearly independent, and the cotangent space has dimension
two.  This exactly matches the Jacobian calculation.  The
agreement is not accidental: the Jacobian method is the
coordinate form of the universal derivation construction.

\paragraph{Relative differentials and the next geometric question.}
For a ring map \(A\to B\), \(\Omega_{B/A}\) records only
changes in \(B\) that hold \(A\) fixed.  For example,
\(B=A[t]\) has \(\Omega_{B/A}\cong B\,dt\).  If a
family of schemes is described by \(\operatorname{Spec}B\to
\operatorname{Spec}A\), the relative differentials capture
tangent directions \emph{along the fibers}.  Their behavior under
base change and at singular points provides a bridge toward
smoothness, deformation theory and the sheaf of differentials.

\paragraph{Worked family: a smooth total space with a singular fiber.}
Return to the fiber-product viewpoint from VI/23--VI/25. Set
\(A=k[t]\) and
\[
 B=A[x,y]/(xy-t),\qquad
 f:\operatorname{Spec}B\longrightarrow\operatorname{Spec}A,
\]
where \(k\) has characteristic zero. Relative to \(A\), the
parameter \(t\) is held fixed. Differentiating the relation gives
\(y\,dx+x\,dy=0\), so
\[
 \Omega_{B/A}\cong
 \frac{B\,dx\oplus B\,dy}{\langle y\,dx+x\,dy\rangle}.
\]
Over \(t=c\), the fiber is \(xy=c\). If \(c\ne0\), neither
coordinate vanishes at any point of the fiber, so the relative
Jacobian row \((y,x)\) has rank one everywhere. Every fiber point
has a one-dimensional tangent space. For \(c=0\), the fiber is
the union of the coordinate axes. At their intersection
\((x,y)=(0,0)\), the Jacobian row vanishes, and the tangent
space has dimension two. The special fiber is therefore singular
even though the total coordinate ring \(B\cong k[x,y]\) is a
regular polynomial ring.

This separates three different notions: regularity of the
\emph{total space}, regularity of a \emph{fiber}, and smoothness
of the \emph{morphism}. Relative differentials make the
appearance of the singular fiber visible in a direct computation.
